% Capitulo: Arquitetura medalhão e modelagem tradicional
\chapter{Arquitetura medalhão e modelagem tradicional}
\label{cap:medalhao}

Esta parte do livro trata de como organizar os dados, agora que eles precisam servir a
muito mais do que painéis. O ponto de partida é uma arquitetura que se popularizou na era
do lakehouse e virou vocabulário comum em times de dados: a arquitetura medalhão, com suas
camadas bronze, prata e ouro. Este capítulo explica o que ela é, em que se parece e em que
difere da modelagem tradicional que o capítulo~\ref{cap:dw} apresentou, e por que ela é um
bom ponto de partida para pensar em vários usos do mesmo dado. A tese é simples: medalhão
organiza camadas de refinamento; não substitui modelagem.

\section{Bronze, prata e ouro}
\label{sec:bronze-prata-ouro}

O nome vem do ecossistema do lakehouse, em que o termo foi difundido para descrever uma
organização dos dados em camadas de qualidade crescente
\parencite{databricksmedallion,armbrust2021}. Cada camada tem um papel claro.

A camada \textbf{bronze} guarda o dado como chegou da fonte, sem transformação. Se o
sistema de origem mandou um arquivo com colunas a mais, datas em formato estranho e
registros duplicados, é assim que ele fica no bronze, junto com o momento em que chegou.
Parece desperdício guardar dado sujo, mas o bronze tem duas funções importantes: permite
reprocessar tudo quando uma regra de negócio muda, e permite auditar, mostrando exatamente
o que a fonte mandou e quando.

A camada \textbf{prata} guarda o dado limpo e integrado. Os tipos estão corretos, as
duplicidades foram tratadas, os nomes de colunas seguem um padrão, e as chaves de sistemas
diferentes foram conciliadas, de forma que o mesmo cliente tem o mesmo identificador
independentemente da fonte. A prata ainda não está organizada para um uso específico; é
uma base confiável a partir da qual vários usos podem ser construídos.

A camada \textbf{ouro} guarda o dado pronto para um consumo específico. Para o BI, modelos
dimensionais e métricas de negócio. Para a ciência de dados, tabelas de variáveis. Para
um assistente de IA, documentos preparados e índices de busca. Pode haver muitas tabelas
ouro construídas sobre a mesma prata, cada uma desenhada para quem vai usá-la.

Quem trabalha com transformação de dados como código vai reconhecer aqui as camadas de
preparação, intermediária e de consumo descritas no capítulo~\ref{cap:nuvem}
(\veja{sec:elt}). A correspondência não é exata, mas a ideia é a mesma: separar o trabalho
de limpar do trabalho de modelar para um uso.

\begin{figure}[htbp]
  \centering
  % Arquitetura medalhao e os consumidores de cada camada (cap. 12).
\begin{tikzpicture}[diagrama, x=1mm, y=1mm]
  % fontes
  \foreach \y/\nome/\id in {12/Sistemas/f1, 4/APIs/f2, -4/Arquivos/f3, -12/Documentos/f4}
    \node[caixa, minimum width=15mm, minimum height=6mm] (\id) at (7.5,\y) {\nome};
  \node[moldura, fit=(f1)(f4), label={[rotulo]above:Fontes}] {};

  % camadas
  \node[banco, minimum width=19mm, minimum height=26mm] (b) at (31,0)
    {\textbf{Bronze}\\[1mm]{\scriptsize bruto, como}\\{\scriptsize chegou}};
  \node[banco, minimum width=19mm, minimum height=26mm] (p) at (53,0)
    {\textbf{Prata}\\[1mm]{\scriptsize limpo, tipado,}\\{\scriptsize integrado}};
  \node[bancodestaque, minimum width=19mm, minimum height=26mm] (o) at (75,0)
    {\textbf{Ouro}\\[1mm]{\scriptsize pronto para}\\{\scriptsize o consumo}};
  \foreach \id in {f1,f2,f3,f4} \draw[seta] (\id.east) -- (b.west);
  \draw[seta] (b) -- (p);
  \draw[seta] (p) -- (o);

  % consumidores
  \foreach \y/\nome/\id in {15/{BI descritivo}/c1, 5/{Ciência de dados}/c2, -5/{Treino de IA}/c3, -15/{Contexto p/ IA}/c4}
    \node[caixa, minimum width=21mm, minimum height=6mm] (\id) at (104.5,\y) {\nome};
  \foreach \id in {c1,c2,c3,c4} \draw[seta] (o.east) -- (\id.west);
  \node[rotulo, text width=110mm] at (58,-23)
    {Ciência de dados e treino de IA muitas vezes leem também da prata, que guarda mais detalhe que o ouro.};
\end{tikzpicture}

  \caption{A arquitetura medalhão e os consumidores que leem de cada camada.}
  \label{fig:medalhao}
  \fonte{elaborado pelo autor com base em \textcite{databricksmedallion}.}
\end{figure}

A Figura~\ref{fig:medalhao} mostra as três camadas e os consumidores típicos. Repare que
nem todo consumo vem do ouro. Ciência de dados e treino de IA muitas vezes leem direto da
prata, porque precisam de detalhe que o ouro, agregado para o BI, já perdeu.

\section{O que muda em relação a Inmon e Kimball}
\label{sec:medalhao-tradicional}

É comum ver a arquitetura medalhão apresentada como alternativa a Inmon e Kimball. Não é.
As duas escolas do capítulo~\ref{cap:dw} falam de \emph{modelo}: como as tabelas devem ser
desenhadas. A arquitetura medalhão fala de \emph{camadas de refinamento}: em que estágio de
tratamento cada dado está. As duas coisas se combinam.

Na prática, a prata lembra muito a camada integrada de Inmon: um lugar onde os dados das
várias fontes são limpos e conciliados antes de servir a qualquer uso. E o ouro, quando
serve ao BI, quase sempre é um modelo dimensional de Kimball, com fatos, dimensões e grão
declarado (\veja{sec:modelagem-dimensional}). O bronze é a novidade mais clara em relação
aos armazéns dos anos 1990, que raramente guardavam o dado bruto, porque armazenamento era
caro. Com armazenamento barato, guardar o bruto deixou de ser luxo e passou a ser
recomendável.

O risco está em tratar a arquitetura medalhão como receita completa. Três camadas não
garantem nada sobre a qualidade do desenho. É possível ter bronze, prata e ouro e, ainda
assim, um ouro sem grão definido, com a mesma métrica calculada de três formas em três
tabelas, e uma prata em que a conciliação de clientes nunca foi feita. Nesse caso, o
resultado é a mesma bagunça de antes, só que organizada em pastas com nomes de metais.

\begin{destaque}
A arquitetura medalhão diz em que estágio cada dado está. A modelagem diz como ele deve ser
desenhado. Uma não substitui a outra.
\end{destaque}

\rascunho{
  Relato: como eu organizei as camadas num projeto real (o que ficou em cada camada, o que
  foi difícil de decidir, que regra de nomenclatura funcionou). Sem nomes de empresa.
}

\section{Modelagem para análise descritiva}
\label{sec:modelagem-descritiva}

A análise descritiva responde ao que aconteceu e por quê. É o uso que a modelagem
tradicional sempre atendeu, e continua sendo, de longe, o maior volume de trabalho na
maioria das empresas: acompanhamento de vendas, custos, operação, metas.

Na camada ouro, a modelagem para análise descritiva segue os princípios do
capítulo~\ref{cap:dw}. Esquemas estrela, com o grão declarado e as dimensões conformadas.
Métricas definidas uma vez, preferencialmente numa camada semântica, e reaproveitadas por
todos os painéis (\veja{sec:camada-semantica}). Tabelas agregadas para os painéis mais
acessados, quando o custo de calcular na hora não compensa.

Algumas práticas da tradição valem especialmente aqui. Chaves substitutas protegem o modelo
de mudanças nas fontes. O histórico das dimensões, guardado com datas de validade, permite
responder como as coisas eram em qualquer momento do passado. E nomes claros, que um
analista novo entende sem perguntar, são talvez o investimento mais barato e mais
subestimado de todos. Eles ajudam o analista novo e, como mostra o
capítulo~\ref{cap:contexto-ia}, ajudam também a IA.

\section{Um dado, vários usos}
\label{sec:um-dado-varios-usos}

A modelagem tradicional tinha um consumidor principal: o analista, usando uma ferramenta de
BI. Hoje, o mesmo dado alimenta modelos preditivos, decisões automatizadas, o treino e o
ajuste de modelos de IA, e assistentes que respondem perguntas em linguagem natural. Cada
um desses usos pede uma forma diferente de organizar os dados, e um dos erros mais comuns é
tentar servir a todos com o mesmo modelo, geralmente o que foi feito para o BI.

A Tabela~\ref{tab:usos-modelagem} resume as diferenças e serve de mapa para os próximos
capítulos desta parte. O capítulo~\ref{cap:master-data} trata da base que todos os usos
compartilham, os cadastros mestres. Os capítulos~\ref{cap:governanca}
e~\ref{cap:aplicacao-propria} tratam de como cuidar dessa base. E os
capítulos~\ref{cap:modelagem-cd}, \ref{cap:treinar-ia} e~\ref{cap:contexto-ia} tratam de
um uso cada um.

\begin{table}[htbp]
  \centering
  \caption{Cada uso pede uma forma diferente de organizar os dados.}
  \label{tab:usos-modelagem}
  \small
  \begin{tabularx}{\textwidth}{@{}>{\raggedright\arraybackslash}p{0.22\textwidth}
      >{\raggedright\arraybackslash}X>{\raggedright\arraybackslash}X@{}}
    \toprule
    Uso & Forma dos dados & Cuidado principal \\
    \midrule
    Descritiva \newline{\footnotesize\itshape o que aconteceu?} & Esquema estrela e tabelas agregadas, na camada ouro & Métrica definida uma vez; grão explícito \\
    \addlinespace
    Preditiva \newline{\footnotesize\itshape o que vai acontecer?} & Tabela larga: uma linha por entidade e momento, com variáveis e alvo & Não deixar o futuro vazar para o passado \\
    \addlinespace
    Prescritiva \newline{\footnotesize\itshape o que fazer?} & Previsões junto com custos, restrições e cenários do negócio & Registrar a decisão tomada e o resultado \\
    \addlinespace
    Treino de IA & Conjunto curado, rotulado, versionado e documentado & Qualidade e representatividade dos exemplos, privacidade \\
    \addlinespace
    Contexto para IA & Documentos em trechos, metadados, glossário, camada semântica & Atualização, dono de cada fonte, permissões \\
    \bottomrule
  \end{tabularx}
  \fonte{elaborado pelo autor.}
\end{table}

A boa notícia é que a arquitetura em camadas facilita essa diversidade. Com uma prata
confiável, construir um ouro novo para um uso novo é muito mais barato do que voltar às
fontes. O investimento que mais rende, portanto, é na prata: limpeza, integração e,
principalmente, a conciliação dos cadastros, que é o assunto do próximo capítulo.

\begin{leituras}
  \leitura{armbrust2021}{O artigo que propõe o lakehouse, o contexto em que a
    arquitetura medalhão se popularizou.}
  \leitura{kleppmann2017}{Para entender os sistemas por baixo das camadas:
    armazenamento, replicação, lote e fluxo contínuo, sem depender de ferramenta.}
  \leitura{kimball2013}{Para a modelagem da camada ouro voltada ao BI, que continua sendo a
    dimensional.}
\end{leituras}
